\documentclass[12pt, a4paper]{article}

% --- Pacotes Essenciais ---
\usepackage[utf8]{inputenc}
\usepackage[T1]{fontenc}
\usepackage[brazil]{babel} % Para formatação em português do Brasil
\usepackage{amsmath}       % Para equações matemáticas avançadas
\usepackage{amssymb}       % Para símbolos matemáticos
\usepackage{graphicx}      % Para incluir imagens
\usepackage{hyperref}      % Para criar links internos e externos (ex: no sumário)
\usepackage{geometry}      % Para ajustar as margens
\usepackage{minted}        % Para código python

% --- Configuração das Margens ---
\geometry{
 a4paper,
 top=3cm,
 bottom=2cm,
 left=3cm,
 right=2cm
}

% --- Informações do Título ---
\title{Grupo 4 - Programação Linear}

% --- Múltiplos Autores ---
% Use \and para separar os autores. Para mais organização, pode-se usar uma tabela.
\author{
  Bruno Ferreira da Silva \
  \small{UERJ} \
  \small{\texttt{email1@exemplo.com}}
  \and
  Maiana Souza Oliveira \
  \small{UERJ} \
  \small{\texttt{email2@exemplo.com}}
  \and
  Marcelo Henrique Araújo Silva \
  \small{UERJ} \
  \small{\texttt{email3@exemplo.com}}
  \and
  Henrique Messias Pimenta \
  \small{UERJ} \
  \small{\texttt{email4@exemplo.com}}
}

\date{\today} % Ou defina uma data específica, ex: \date{Setembro de 2023}

% --- Início do Documento ---
\begin{document}

\maketitle % Cria a página de título com as informações acima
\cleardoublepage

\begin{abstract}
    \noindent
    Este trabalho tem como objetivo demonstrar como os alunos desenvolveram ao longo do curso suas capacidades de resolver problemas de programação linear com base em dois problemas: Um exercício de aplicação dos métodos simplex e um problema de modelagem.
    \vspace{\baselineskip} % Adiciona um espaço vertical
    
    \noindent
    \textbf{Palavras-chave:} Programação Linear, Otimização, Método Simplex, [Outra Palavra-chave].
\end{abstract}

\newpage
\tableofcontents % Gera o sumário automagicamente
\newpage

% --- Seções do Projeto ---

\section{Introdução ao método simplex}
[Resumo de problemas do programação linear]
[Resumo do método simplex analítico]
[Resumo do método simplex tabular]

\subsection{Formulação de Problemas}
[Explique como modelar um problema real como um problema de programação linear (PPL). Se possível, use um exemplo clássico, como o problema da dieta ou da alocação de recursos.]

\subsubsection{Problema 1}
Agora vamos usar tudo o que sabemos para resolver o seguinte problema:
\begin{equation*}
    \min Z = 2x_1 + x_2, \qquad\text{s, r} \left\{\begin{matrix}
         x_1 +  x_2 \ge 3 \\
        2x_1 + 3x_2 \ge 9 \\
         x_1        \ge 1 \\
         x_1,  x_2  \ge 0
    \end{matrix}\right.
\end{equation*}
[Descreva os principais algoritmos para resolver PPLs, como o método gráfico (para duas variáveis) e, principalmente, o Método Simplex.]

\subsubsection{Problema 2}
A Companhia Coelho S.A. fabrica motores para brinquedos e pequenos aparelhos. O departamento de marketing está prevendo vendas de 6.100 unidades do motor Roncam no próximo semestre. Essa é uma nova demanda, e a companhia terá de testar sua capacidade produtiva. O motor Roncam é montado com três componentes: um corpo, uma base e uma blindagem. Alguns desses componentes podem ser comprados de outros fornecedores, se existirem limitações da Coelho S.A. em sua produção. Os custos de produção e de aquisição em R\$/unidade estão resumidos na tabela a seguir:

\begin{table*}[htb]
  \begin{tabular}{ccc}
    Componentes & Custo de aquisição (R\$) & Custos de produção (R\$) \\
    Corpo & 10,00 & 8,00 \\
    Base & 20,00 & 20,00 \\
    Blindagem & 20,00 & 10,00
  \end{tabular}
\end{table*}

A fábrica da Companhia Coelho S.A. tem três departamentos. O requisito de tempo em horas que cada componente consome em cada departamento está resumido na tabela a seguir. As horas disponíveis para a companhia estão listadas na última linha.

\begin{table*}[htb]
  \begin{tabular}{cccc}
    Componentes & Preparação & Molde & Fabricação\\
    Corpo & 1/30 & 1/15 & 1/30 \\
    Base & 1/12 & 1/30 & 1/15 \\
    Blindagem & 1/15 & 1/15 & 1/12 \\
    Disponibilidade & 820 & 820 & 820 \\
  \end{tabular}
\end{table*}

\section{Metodologia}
[Descreva os passos que você seguiu para desenvolver o projeto. Isso pode incluir a ferramenta ou linguagem de programação utilizada (ex: Python com a biblioteca SciPy, MATLAB, LINGO), a fonte dos dados (se aplicável) e o design do experimento ou da aplicação.]

\section{Desenvolvimento e Resultados}
[Apresente o problema específico que você resolveu. Mostre a formulação matemática, o código ou a implementação utilizada e os resultados obtidos. Use tabelas, gráficos e figuras para ilustrar suas descobertas.]

% Exemplo de como incluir uma figura
% \begin{figure}[h!]
%     \centering
%     \includegraphics[width=0.8\textwidth]{caminho/para/sua/imagem.png}
%     \caption{Descrição da sua figura.}
%     \label{fig:exemplo}
% \end{figure}

\section{Conclusão}
% --- Seção de Agradecimentos (não numerada) ---
\section*{Agradecimentos}
\addcontentsline{toc}{section}{Agradecimentos} % Adiciona ao sumário
[Este é o espaço para agradecer aos seus orientadores, colegas, instituições de fomento (como CAPES, CNPq), amigos e familiares que contribuíram de alguma forma para a realização do projeto.]

% --- Referências Bibliográficas ---
\begin{thebibliography}{99}
\addcontentsline{toc}{section}{Referências} % Adiciona ao sumário

\bibitem{exemplo_livro}
SOBRENOME, N. \textit{Título do Livro}. Edição. Cidade: Editora, Ano.

\bibitem{exemplo_artigo}
SOBRENOME, N.; SOBRENOME, N. Título do Artigo. \textit{Nome do Periódico}, v. X, n. Y, p. 1-10, Ano.

% Adicione suas referências aqui...

\end{thebibliography}

\end{document}
